\subsection{结果与机制解释}
表 \ref{tab:q1_result} 给出了 2025 年的总量标定结果与局地对照结果。四大家鱼总量已较好对齐题目给出的恢复信息；$\mathrm{CPUE}^*$ 虽低于 2.2 的湖北段对照值，但恢复方向一致。
\begin{table}[htbp]
\centering
\caption{问题一的 2025 年参数标定结果与局地对照}
\label{tab:q1_result}
\begin{tabularx}{0.95\textwidth}{>{\raggedright\arraybackslash}p{3.0cm} >{\raggedright\arraybackslash}p{2.6cm} >{\raggedright\arraybackslash}p{2.6cm} >{\raggedright\arraybackslash}p{2.8cm}X}
\toprule
指标 & 观测值 & 模拟值 & 相对误差 & 说明 \\
\midrule
四大家鱼总量倍率 & 1.800 & 1.805 & 0.26\% & 对齐 \\
$\mathrm{CPUE}^*$ 倍率 & 2.200 & 1.838 & 16.47\% & 一致 \\
\bottomrule
\end{tabularx}
\end{table}

这两项指标变化方向一致，但恢复幅度存在差异。尽管统计口径不同，二者均证实了禁渔后的资源回升趋势：四大家鱼总生物量与1.8倍的标定值几乎完全吻合，表明单区模型能够准确刻画全流域平均恢复水平；$\mathrm{CPUE}^*$低于湖北江段实测的2.2倍，则反映出局地高捕捞效率条件下资源增幅可能更为显著。

从动态过程来看，禁渔情景下四大家鱼总量与$\mathrm{CPUE}^*$同步上升，但基础资源末值仅为初始值的0.204倍，且水草生物量最小值出现在禁渔第10年，说明鱼类资源恢复并未随基础资源状况一同改善。摄食压力指数从初始水平持续攀升至4.213，表明鱼群恢复程度与基础资源承压强度呈正相关。0.204与$\Phi$均为归一化单区系统中的相对指标，不直接对应全流域实际生物量或能量金字塔的绝对比例。$\Phi$的持续上升与基础资源的大幅下降共同表明，资源供给已无法满足消费者种群的快速增长；禁渔解除捕捞压力后，中层鱼类对基础资源的摄食压力与种间竞争显著增强，系统呈现出"鱼类资源恢复与基础资源承压并存"的特征。

\begin{figure}[htbp]
\centering
\includegraphics[width=0.92\textwidth]{fig01_q1_foodweb.png}
\caption{问题一：标定后的禁渔与对照情景及观测约束}
\label{fig:q1}
\end{figure}

